\documentclass[12pt, a4paper]{article}

% ==================== 宏包引入 ====================
\usepackage[UTF8]{ctex}           % 中文支持
\usepackage{geometry}             % 页面布局
\usepackage{amsmath, amssymb}     % 数学公式
\usepackage{graphicx}             % 插图
\usepackage{hyperref}             % 超链接
\usepackage{listings}             % 代码高亮
\usepackage{xcolor}               % 颜色
\usepackage{booktabs}             % 三线表
\usepackage{enumitem}             % 列表定制
\usepackage{fancyhdr}             % 页眉页脚

\geometry{left=2.5cm, right=2.5cm, top=2.5cm, bottom=2.5cm}

% 代码样式
\lstset{
  basicstyle=\small\ttfamily,
  keywordstyle=\color{blue}\bfseries,
  commentstyle=\color{gray},
  stringstyle=\color{red},
  backgroundcolor=\color{gray!10},
  frame=single,
  breaklines=true,
  language=[LaTeX]TeX,
  literate={\\}{{{\color{blue}\textbackslash}}}1,
}

\hypersetup{colorlinks=true, linkcolor=blue, urlcolor=cyan}

\title{\LaTeX{} 入门教程}
\author{自动生成}
\date{\today}

\begin{document}

\maketitle
\tableofcontents
\newpage

% ==================== 第一章 ====================
\section{什么是 \LaTeX{}}

\LaTeX{} 是一种基于 \TeX{} 的排版系统，由 Leslie Lamport 开发。它特别适合：

\begin{itemize}
  \item 学术论文、书籍排版
  \item 数学公式的精确呈现
  \item 自动化的交叉引用、目录和参考文献管理
\end{itemize}

与 Word 等所见即所得编辑器不同，\LaTeX{} 采用"所写即所得"的方式——你编写源代码，编译后生成 PDF。

% ==================== 第二章 ====================
\section{文档结构}

一个最简单的 \LaTeX{} 文档如下：

\begin{lstlisting}
\documentclass{article}
\begin{document}
Hello, World!
\end{document}
\end{lstlisting}

\subsection{文档类}

\verb|\documentclass| 指定文档类型，常用选项：

\begin{table}[h]
\centering
\begin{tabular}{ll}
\toprule
文档类 & 用途 \\
\midrule
\texttt{article} & 短文、论文 \\
\texttt{report}  & 长篇报告（含章节） \\
\texttt{book}    & 书籍 \\
\texttt{beamer}  & 幻灯片 \\
\bottomrule
\end{tabular}
\caption{常用文档类}
\end{table}

\subsection{前言区与正文区}

\begin{lstlisting}
\documentclass[12pt, a4paper]{article}
% ----- 前言区 (preamble) -----
\usepackage{amsmath}
\title{我的文档}
\author{张三}

% ----- 正文区 -----
\begin{document}
\maketitle
正文内容...
\end{document}
\end{lstlisting}

前言区用于加载宏包、设置参数；正文区包含实际内容。

% ==================== 第三章 ====================
\section{文本排版}

\subsection{字体样式}

\begin{tabular}{ll}
\toprule
命令 & 效果 \\
\midrule
\verb|\textbf{粗体}| & \textbf{粗体} \\
\verb|\textit{斜体}| & \textit{斜体} \\
\verb|\underline{下划线}| & \underline{下划线} \\
\verb|\texttt{等宽}| & \texttt{等宽} \\
\bottomrule
\end{tabular}

\subsection{段落与换行}

\begin{itemize}
  \item 空行产生新段落
  \item \verb|\\| 强制换行（不产生新段落）
  \item \verb|\newpage| 强制换页
\end{itemize}

\subsection{列表}

\textbf{无序列表：}
\begin{lstlisting}
\begin{itemize}
  \item 第一项
  \item 第二项
\end{itemize}
\end{lstlisting}

\textbf{有序列表：}
\begin{lstlisting}
\begin{enumerate}
  \item 第一项
  \item 第二项
\end{enumerate}
\end{lstlisting}

% ==================== 第四章 ====================
\section{数学公式}

这是 \LaTeX{} 最强大的能力之一。

\subsection{行内公式}

用 \verb|$...$| 包裹：质能方程 $E = mc^2$。

\subsection{行间公式}

用 \verb|\[...\]| 或 \texttt{equation} 环境：

\begin{equation}
  \int_{-\infty}^{\infty} e^{-x^2} \, dx = \sqrt{\pi}
\end{equation}

\subsection{常用数学符号}

\begin{table}[h]
\centering
\begin{tabular}{lll}
\toprule
代码 & 效果 & 说明 \\
\midrule
\verb|\frac{a}{b}| & $\frac{a}{b}$ & 分数 \\
\verb|x^{2}| & $x^{2}$ & 上标 \\
\verb|x_{i}| & $x_{i}$ & 下标 \\
\verb|\sqrt{x}| & $\sqrt{x}$ & 根号 \\
\verb|\sum_{i=1}^{n}| & $\sum_{i=1}^{n}$ & 求和 \\
\verb|\alpha, \beta| & $\alpha, \beta$ & 希腊字母 \\
\verb|\leq, \geq| & $\leq, \geq$ & 不等号 \\
\verb|\infty| & $\infty$ & 无穷 \\
\bottomrule
\end{tabular}
\caption{常用数学符号速查}
\end{table}

\subsection{矩阵}

\begin{lstlisting}
\begin{equation}
  A = \begin{pmatrix}
    1 & 2 \\
    3 & 4
  \end{pmatrix}
\end{equation}
\end{lstlisting}

效果：
\begin{equation}
  A = \begin{pmatrix}
    1 & 2 \\
    3 & 4
  \end{pmatrix}
\end{equation}

\subsection{多行公式对齐}

\begin{lstlisting}
\begin{align}
  f(x) &= x^2 + 2x + 1 \\
       &= (x + 1)^2
\end{align}
\end{lstlisting}

效果：
\begin{align}
  f(x) &= x^2 + 2x + 1 \\
       &= (x + 1)^2
\end{align}

% ==================== 第五章 ====================
\section{图片与表格}

\subsection{插入图片}

\begin{lstlisting}
\begin{figure}[htbp]
  \centering
  \includegraphics[width=0.6\textwidth]{example.png}
  \caption{示例图片}
  \label{fig:example}
\end{figure}
\end{lstlisting}

浮动位置参数：\texttt{h}(当前位置)、\texttt{t}(页顶)、\texttt{b}(页底)、\texttt{p}(独立页)。

\subsection{三线表}

\begin{lstlisting}
\begin{table}[htbp]
  \centering
  \begin{tabular}{ccc}
    \toprule
    姓名 & 年龄 & 成绩 \\
    \midrule
    张三 & 20 & 95 \\
    李四 & 21 & 88 \\
    \bottomrule
  \end{tabular}
  \caption{学生成绩表}
\end{table}
\end{lstlisting}

% ==================== 第六章 ====================
\section{交叉引用与参考文献}

\subsection{标签与引用}

\begin{lstlisting}
\section{方法}\label{sec:method}
如公式 \ref{eq:euler} 所示...

\begin{equation}\label{eq:euler}
  e^{i\pi} + 1 = 0
\end{equation}
\end{lstlisting}

使用 \verb|\label{}| 打标签，\verb|\ref{}| 引用编号，\verb|\pageref{}| 引用页码。

\subsection{参考文献（BibTeX）}

创建 \texttt{refs.bib} 文件：

\begin{lstlisting}
@article{einstein1905,
  author  = {Albert Einstein},
  title   = {On the Electrodynamics of Moving Bodies},
  journal = {Annalen der Physik},
  year    = {1905},
}
\end{lstlisting}

在文档中引用：

\begin{lstlisting}
根据 Einstein \cite{einstein1905} 的研究...

\bibliographystyle{plain}
\bibliography{refs}
\end{lstlisting}

% ==================== 第七章 ====================
\section{常用宏包速查}

\begin{table}[h]
\centering
\begin{tabular}{ll}
\toprule
宏包 & 功能 \\
\midrule
\texttt{ctex} & 中文排版支持 \\
\texttt{amsmath} & 增强数学公式 \\
\texttt{graphicx} & 插入图片 \\
\texttt{hyperref} & 超链接与 PDF 书签 \\
\texttt{listings} & 代码高亮 \\
\texttt{geometry} & 自定义页面边距 \\
\texttt{booktabs} & 专业三线表 \\
\texttt{tikz} & 矢量绘图 \\
\texttt{biblatex} & 现代参考文献管理 \\
\bottomrule
\end{tabular}
\caption{常用宏包}
\end{table}

% ==================== 第八章 ====================
\section{编译方式}

\subsection{常用编译命令}

\begin{lstlisting}[language=bash]
# 推荐：XeLaTeX（支持中文和 Unicode）
xelatex tutorial.tex

# 含参考文献时的完整流程
xelatex tutorial.tex
bibtex tutorial
xelatex tutorial.tex
xelatex tutorial.tex

# 或使用 latexmk 自动处理
latexmk -xelatex tutorial.tex
\end{lstlisting}

\subsection{推荐编辑器}

\begin{itemize}
  \item \textbf{VS Code} + LaTeX Workshop 插件
  \item \textbf{Overleaf}（在线协作）
  \item \textbf{TeXstudio}（专用 IDE）
\end{itemize}

\section{总结}

本教程涵盖了 \LaTeX{} 的核心用法：文档结构、文本排版、数学公式、图表、交叉引用和参考文献。掌握这些内容足以应对大部分学术写作需求。更多进阶内容可参考：

\begin{itemize}
  \item \href{https://www.overleaf.com/learn}{Overleaf 官方文档}
  \item \href{https://ctan.org/}{CTAN 宏包仓库}
  \item 《\LaTeX 入门》（刘海洋 著）
\end{itemize}

\end{document}
